% 原创TikZ重绘；层次布局参考Bishop与Bishop（2024）第12章。
\begin{tikzpicture}[x=1mm,y=1mm,font=\figurefont\footnotesize,
 every node/.style={text=NNInk},
 stage/.style={nn operator,minimum height=9mm,inner sep=2pt},
 norm/.style={stage,minimum height=7mm},
 wire/.style={nn flow,draw=NNWire,rounded corners=1.5mm},
 aux/.style={draw=NNLine,line width=.85pt,densely dashed},
 sum/.style={nn pointwise,minimum size=5mm}]

\node[] (gtok0) at (28,0) {猫};
\node[inner sep=0pt,minimum width=12mm,minimum height=3mm] (ge0) at (28,12) {};
\filldraw[draw=NNInput,fill=NNInput!18,line width=.5pt] (22.0,10.5) rectangle ++(3.0,3);
\filldraw[draw=NNInput,fill=NNInput!18,line width=.5pt] (25.0,10.5) rectangle ++(3.0,3);
\filldraw[draw=NNInput,fill=NNInput!18,line width=.5pt] (28.0,10.5) rectangle ++(3.0,3);
\filldraw[draw=NNInput,fill=NNInput!18,line width=.5pt] (31.0,10.5) rectangle ++(3.0,3);
\draw[wire] (gtok0.north)--(ge0.south);
\node[] (gplus0) at (28,18) {$+$};
\node[inner sep=0pt,minimum width=12mm,minimum height=3mm] (gpos0) at (28,24) {};
\filldraw[draw=NNInput,fill=NNInput!18,line width=.5pt] (22.0,22.5) rectangle ++(3.0,3);
\filldraw[draw=NNInput,fill=NNInput!18,line width=.5pt] (25.0,22.5) rectangle ++(3.0,3);
\filldraw[draw=NNInput,fill=NNInput!18,line width=.5pt] (28.0,22.5) rectangle ++(3.0,3);
\filldraw[draw=NNInput,fill=NNInput!18,line width=.5pt] (31.0,22.5) rectangle ++(3.0,3);
\draw[aux] (21,8) rectangle (35,28);
\node[inner sep=0pt,minimum width=12mm,minimum height=3mm] (gh00) at (28,39) {};
\filldraw[draw=NNInput,fill=NNInput!18,line width=.5pt] (22.0,37.5) rectangle ++(3.0,3);
\filldraw[draw=NNInput,fill=NNInput!18,line width=.5pt] (25.0,37.5) rectangle ++(3.0,3);
\filldraw[draw=NNInput,fill=NNInput!18,line width=.5pt] (28.0,37.5) rectangle ++(3.0,3);
\filldraw[draw=NNInput,fill=NNInput!18,line width=.5pt] (31.0,37.5) rectangle ++(3.0,3);
\draw[wire] (28,28)--(gh00.south);
\node[] (gtok1) at (64,0) {追};
\node[inner sep=0pt,minimum width=12mm,minimum height=3mm] (ge1) at (64,12) {};
\filldraw[draw=NNInput,fill=NNInput!18,line width=.5pt] (58.0,10.5) rectangle ++(3.0,3);
\filldraw[draw=NNInput,fill=NNInput!18,line width=.5pt] (61.0,10.5) rectangle ++(3.0,3);
\filldraw[draw=NNInput,fill=NNInput!18,line width=.5pt] (64.0,10.5) rectangle ++(3.0,3);
\filldraw[draw=NNInput,fill=NNInput!18,line width=.5pt] (67.0,10.5) rectangle ++(3.0,3);
\draw[wire] (gtok1.north)--(ge1.south);
\node[] (gplus1) at (64,18) {$+$};
\node[inner sep=0pt,minimum width=12mm,minimum height=3mm] (gpos1) at (64,24) {};
\filldraw[draw=NNInput,fill=NNInput!18,line width=.5pt] (58.0,22.5) rectangle ++(3.0,3);
\filldraw[draw=NNInput,fill=NNInput!18,line width=.5pt] (61.0,22.5) rectangle ++(3.0,3);
\filldraw[draw=NNInput,fill=NNInput!18,line width=.5pt] (64.0,22.5) rectangle ++(3.0,3);
\filldraw[draw=NNInput,fill=NNInput!18,line width=.5pt] (67.0,22.5) rectangle ++(3.0,3);
\draw[aux] (57,8) rectangle (71,28);
\node[inner sep=0pt,minimum width=12mm,minimum height=3mm] (gh01) at (64,39) {};
\filldraw[draw=NNInput,fill=NNInput!18,line width=.5pt] (58.0,37.5) rectangle ++(3.0,3);
\filldraw[draw=NNInput,fill=NNInput!18,line width=.5pt] (61.0,37.5) rectangle ++(3.0,3);
\filldraw[draw=NNInput,fill=NNInput!18,line width=.5pt] (64.0,37.5) rectangle ++(3.0,3);
\filldraw[draw=NNInput,fill=NNInput!18,line width=.5pt] (67.0,37.5) rectangle ++(3.0,3);
\draw[wire] (64,28)--(gh01.south);
\node[] (gtok2) at (100,0) {狗};
\node[inner sep=0pt,minimum width=12mm,minimum height=3mm] (ge2) at (100,12) {};
\filldraw[draw=NNInput,fill=NNInput!18,line width=.5pt] (94.0,10.5) rectangle ++(3.0,3);
\filldraw[draw=NNInput,fill=NNInput!18,line width=.5pt] (97.0,10.5) rectangle ++(3.0,3);
\filldraw[draw=NNInput,fill=NNInput!18,line width=.5pt] (100.0,10.5) rectangle ++(3.0,3);
\filldraw[draw=NNInput,fill=NNInput!18,line width=.5pt] (103.0,10.5) rectangle ++(3.0,3);
\draw[wire] (gtok2.north)--(ge2.south);
\node[] (gplus2) at (100,18) {$+$};
\node[inner sep=0pt,minimum width=12mm,minimum height=3mm] (gpos2) at (100,24) {};
\filldraw[draw=NNInput,fill=NNInput!18,line width=.5pt] (94.0,22.5) rectangle ++(3.0,3);
\filldraw[draw=NNInput,fill=NNInput!18,line width=.5pt] (97.0,22.5) rectangle ++(3.0,3);
\filldraw[draw=NNInput,fill=NNInput!18,line width=.5pt] (100.0,22.5) rectangle ++(3.0,3);
\filldraw[draw=NNInput,fill=NNInput!18,line width=.5pt] (103.0,22.5) rectangle ++(3.0,3);
\draw[aux] (93,8) rectangle (107,28);
\node[inner sep=0pt,minimum width=12mm,minimum height=3mm] (gh02) at (100,39) {};
\filldraw[draw=NNInput,fill=NNInput!18,line width=.5pt] (94.0,37.5) rectangle ++(3.0,3);
\filldraw[draw=NNInput,fill=NNInput!18,line width=.5pt] (97.0,37.5) rectangle ++(3.0,3);
\filldraw[draw=NNInput,fill=NNInput!18,line width=.5pt] (100.0,37.5) rectangle ++(3.0,3);
\filldraw[draw=NNInput,fill=NNInput!18,line width=.5pt] (103.0,37.5) rectangle ++(3.0,3);
\draw[wire] (100,28)--(gh02.south);
\node[stage,minimum width=89mm,minimum height=10mm] (glayer1) at (64.0,57) {第1层：因果自注意力＋FFN};
\node[stage,minimum width=89mm,minimum height=10mm] (glayerL) at (64.0,85) {第$L$层：因果自注意力＋FFN};
\draw[wire] (gh00.north)--(28,52);
\draw[wire] (28,62)--(28,67);
\node[] (gellipsis0) at (28,71) {$\vdots$};
\draw[wire] (28,75)--(28,80);
\node[inner sep=0pt,minimum width=12mm,minimum height=3mm] (ghL0) at (28,102) {};
\filldraw[draw=NNOutput,fill=NNOutput!18,line width=.5pt] (22.0,100.5) rectangle ++(3.0,3);
\filldraw[draw=NNOutput,fill=NNOutput!18,line width=.5pt] (25.0,100.5) rectangle ++(3.0,3);
\filldraw[draw=NNOutput,fill=NNOutput!18,line width=.5pt] (28.0,100.5) rectangle ++(3.0,3);
\filldraw[draw=NNOutput,fill=NNOutput!18,line width=.5pt] (31.0,100.5) rectangle ++(3.0,3);
\draw[wire] (28,90)--(ghL0.south);
\draw[wire] (gh01.north)--(64,52);
\draw[wire] (64,62)--(64,67);
\node[] (gellipsis1) at (64,71) {$\vdots$};
\draw[wire] (64,75)--(64,80);
\node[inner sep=0pt,minimum width=12mm,minimum height=3mm] (ghL1) at (64,102) {};
\filldraw[draw=NNOutput,fill=NNOutput!18,line width=.5pt] (58.0,100.5) rectangle ++(3.0,3);
\filldraw[draw=NNOutput,fill=NNOutput!18,line width=.5pt] (61.0,100.5) rectangle ++(3.0,3);
\filldraw[draw=NNOutput,fill=NNOutput!18,line width=.5pt] (64.0,100.5) rectangle ++(3.0,3);
\filldraw[draw=NNOutput,fill=NNOutput!18,line width=.5pt] (67.0,100.5) rectangle ++(3.0,3);
\draw[wire] (64,90)--(ghL1.south);
\draw[wire] (gh02.north)--(100,52);
\draw[wire] (100,62)--(100,67);
\node[] (gellipsis2) at (100,71) {$\vdots$};
\draw[wire] (100,75)--(100,80);
\node[inner sep=0pt,minimum width=12mm,minimum height=3mm] (ghL2) at (100,102) {};
\filldraw[draw=NNOutput,fill=NNOutput!18,line width=.5pt] (94.0,100.5) rectangle ++(3.0,3);
\filldraw[draw=NNOutput,fill=NNOutput!18,line width=.5pt] (97.0,100.5) rectangle ++(3.0,3);
\filldraw[draw=NNOutput,fill=NNOutput!18,line width=.5pt] (100.0,100.5) rectangle ++(3.0,3);
\filldraw[draw=NNOutput,fill=NNOutput!18,line width=.5pt] (103.0,100.5) rectangle ++(3.0,3);
\draw[wire] (100,90)--(ghL2.south);
\node[stage,draw=NNOutput,fill=NNOutput!18,minimum width=24mm,minimum height=8mm] (ghead0) at (28,118) {LM头};
\draw[wire] (ghL0.north)--(ghead0.south);
\node[] (gtarget0) at (28,132) {追};
\draw[wire] (ghead0.north)--(gtarget0.south);
\node[stage,draw=NNOutput,fill=NNOutput!18,minimum width=24mm,minimum height=8mm] (ghead1) at (64,118) {LM头};
\draw[wire] (ghL1.north)--(ghead1.south);
\node[] (gtarget1) at (64,132) {狗};
\draw[wire] (ghead1.north)--(gtarget1.south);
\node[stage,draw=NNOutput,fill=NNOutput!18,minimum width=24mm,minimum height=8mm] (ghead2) at (100,118) {LM头};
\draw[wire] (ghL2.north)--(ghead2.south);
\node[] (gtarget2) at (100,132) {EOS};
\draw[wire] (ghead2.north)--(gtarget2.south);
\draw[aux,draw=NNGradient,rounded corners=1.2mm]
  (ge0.west)--(12,12)--(12,118)--(ghead0.west);
\node[rotate=90,text=FigureInk,fill=none,inner sep=1pt] at (8,65)
  {LM头共享输入词嵌入的转置权重};
\node[nn heading] (gdetailtitle) at (145,138) {GPT-1：Post-LN};
\node[] (gdin) at (145,17) {$\bm H^{(l-1)}$};
\node[stage,minimum width=34mm] (gd0) at (145,34) {因果多头注意力};
\draw[wire] (gdin.north)--(gd0.south);
\node[sum] (gd1) at (145,49) {$+$};
\draw[wire] (gd0.north)--(gd1.south);
\node[norm,minimum width=34mm] (gd2) at (145,62) {LN};
\draw[wire] (gd1.north)--(gd2.south);
\node[stage,minimum width=34mm] (gd3) at (145,82) {逐位置FFN};
\draw[wire] (gd2.north)--(gd3.south);
\node[sum] (gd4) at (145,97) {$+$};
\draw[wire] (gd3.north)--(gd4.south);
\node[norm,minimum width=34mm] (gd5) at (145,110) {LN};
\draw[wire] (gd4.north)--(gd5.south);
\node[] (gdout) at (145,124) {$\bm H^{(l)}$};
\draw[wire] (gd5.north)--(gdout.south);
\draw[wire] (145,24)--(167.0,24)--(167.0,49)--(gd1.east);
\fill[NNWire] (145,24) circle (.55);
\draw[wire] (145,70)--(167.0,70)--(167.0,97)--(gd4.east);
\fill[NNWire] (145,70) circle (.55);
\end{tikzpicture}
